\chapter{Switching Lines}
\label{chap:chapter4}
\index{line switching}
\index{cable!energization}
\index{Petersen coil}

\section{General}
\label{sec:section41}

This chapter examines the energization of unloaded and loaded cable lines. It
also considers the energization of a second line connected to an already
energized busbar and the energization of a cable terminated by a transformer.
The examples illustrate the influence of switching instant, line capacitance,
network inductance, load damping, and transformer residual flux on the
resulting transients.

The switching transients are oscillations at the natural frequencies of the
electrical network. Such oscillations occur whenever the system state changes
abruptly, for example during line energization or at the inception of an earth
fault or short circuit.

\section{Circuit Design}
\label{sec:section42}

Open the file: \texttt{switching\_lines.asc}. The LTspice schematic contains four circuits. Each line is a
\SI{5}{\kilo\metre}, \SI{110}{\kilo\volt} oil-filled cable with a conductor
cross-section of \SI{500}{\milli\metre\squared}. The circuits represent:

\begin{enumerate}
\item energization of an unloaded cable,
\item energization of a cable supplying a \SI{100}{\mega\watt} load,
\item energization of a second unloaded cable connected to an energized busbar,
\item energization of a cable terminated by a transformer with unfavourable
      full residual flux.
\end{enumerate}

Each single line is earth-fault compensated by a \SI{125}{\ampere} Petersen
coil. The double-line arrangement initially uses a compensation setting of \SI{125}{\ampere}, which is increased to\SI{250}{\ampere} at 400 ms. 


\subsection{Adjustment of the Petersen Coil}
\label{sec:petersen-adjustment}

To determine the required compensation, temporarily reduce the load in the
second circuit from \texttt{100e6} to \texttt{100e-6} and run the simulation.
The phase-1 charging-current peak is approximately \SI{60}{\ampere}. For a
shielded single-core cable whose dominant capacitances are phase-to-earth, the
earth-fault charging current is approximately

\begin{equation}
I_{\mathrm{E,peak}}\approx 3I_{0,\mathrm{peak}}
 =3\cdot\SI{60}{\ampere}=\SI{180}{\ampere},
\end{equation}

or approximately \SI{127}{\ampere} RMS. A \SI{125}{\ampere} Petersen coil is
therefore appropriate. After this adjustment, restore the load to
\texttt{100e6}.

\begin{remark}{Applicability of the threefold charging current}{remark41}
The relation $I_{\mathrm E}\approx3I_0$ applies to the zero-sequence charging
current of the cable. Overhead lines have appreciable phase-to-phase
capacitances, so the same simplified interpretation must not be transferred
without checking the complete capacitance matrix.
\end{remark}

\section{Energization of an Unloaded \SI{110}{\kilo\volt} Cable}
\label{sec:section43}
NOTE: Simulate \texttt{Switching\_lines.asc} and choose after simulation \texttt{line1.plt}.
This is done by activating the waveform-viewer window and click in the PlotSettings the menu item "open plot settings file".
 
For the selected switching instant, the Phase A switch closes near a voltage zero crossing. Its peak current is therefore the smallest, approximately
\SI{130}{\ampere}. Immediately after closing, the cable capacitances behave
like short circuits and draw large, rapidly changing charging currents. The
phase-2 and phase-3 current peaks reach approximately \SI{1.5}{\kilo\ampere}.
The oscillations have substantially decayed after approximately
\SI{200}{\milli\second}.

The Phase A voltage transient is likewise the smallest. The Phase B and
Phase C voltages reach approximately \SI{-160}{\kilo\volt} peak, about
\SI{30}{\percent} above the nominal phase-voltage peak. The precise current
and voltage maxima depend strongly on the switching instant. This can be
investigated by changing the switch parameter \texttt{t\_start} from
\SI{40}{\milli\second} to \SI{45}{\milli\second}.

\begin{important}{Origin of the voltage transient}{important41}
The capacitive charging current flows through the inductive source and line
impedances. The resulting voltage drop is superimposed on the source voltage
at the cable terminal. Consequently, the phases with the largest charging
current also show the largest voltage transient.
\end{important}

\subsection{Natural Frequency}
\label{sec:unloaded-line-natural-frequency}

The natural frequency can be identified in LTspice with
\textsc{View}/\textsc{FFT}. For example, an FFT of the Phase B voltage 
shows peaks at \SI{50}{\hertz} and approximately
\SI{1.44}{\kilo\hertz}.

The short-circuit power of the network is \SI{7.5}{\giga\volt\ampere}. Its
nominal reactance is

\begin{calculation}{Natural frequency of the \SI{5}{\kilo\metre} cable}{calculation41}
\begin{align}
X_{\mathrm N}
  &=\frac{U_{\mathrm N}^2}{S_{\mathrm k}}
    =\frac{(\SI{110}{\kilo\volt})^2}
           {\SI{7.5}{\giga\volt\ampere}}
    =\SI{1.613}{\ohm},\\
X_{\mathrm{line}}
  &=\SI{0.11}{\ohm\per\kilo\metre}\cdot\SI{5}{\kilo\metre}
    =\SI{0.55}{\ohm}.
\end{align}

With $\omega=2\pi\cdot\SI{50}{\hertz}$ and
$C'=\SI{400}{\nano\farad\per\kilo\metre}$, a lumped approximation gives

\begin{align}
L&=\frac{X_{\mathrm N}+X_{\mathrm{line}}}{\omega}
  \approx\SI{6.9}{\milli\henry},\\
C&=C'l=\SI{2}{\micro\farad},\\
f_{\mathrm r}&=\frac{1}{2\pi\sqrt{LC}}
  \approx\SI{1.35}{\kilo\hertz}.
\end{align}

The cable is implemented as a lumped T-section. During the initial charging
process, only the source-side half of its series inductance is effective in the
relevant resonant loop. Therefore,

\begin{align}
L_{\mathrm{eff}}
 &=\frac{X_{\mathrm N}+X_{\mathrm{line}}/2}{\omega}
   \approx\SI{6.0}{\milli\henry},\\
f_{\mathrm r,T}
 &=\frac{1}{2\pi\sqrt{L_{\mathrm{eff}}C}}
   \approx\boxed{\SI{1.45}{\kilo\hertz}},
\end{align}

which agrees closely with the simulated \SI{1.44}{\kilo\hertz}.
\end{calculation}

The approximate oscillation frequency can also be read directly in the waveform viewer:
zoom into two adjacent transient peaks and use the time difference and
frequency displayed in the lower-left status line.

\subsection{Petersen-Coil Oscillation}
\label{sec:petersen-oscillation}

The Petersen-coil voltage and current do not immediately reach steady state.
A slow beat of approximately \SI{1}{\hertz} is visible because the compensated
zero-sequence resonant frequency differs slightly from the network frequency.
After the transient has decayed, the residual coil current and beat are nearly
zero.

\section{Energization of a Loaded Line}
\label{sec:Section44}
NOTE: Simulate \texttt{Switching\_lines.asc} and choose after simulation \texttt{line2.plt}.
This is done by activating the waveform-viewer window and click in the PlotSettings the menu item "open plot settings file".
  
When the cable supplies the load, its voltage and current transients are more
strongly damped. The active part of the load absorbs oscillatory energy and
reduces the duration of the switching transient.

\section{Energization of a Second Line}
\label{sec:section45}
NOTE: Simulate \texttt{Switching\_lines.asc} and choose after simulation \texttt{line3.plt}.
This is done by activating the waveform-viewer window and click in the PlotSettings the menu item "open plot settings file".
  
The first, upper line is energized in the same way as the unloaded line in
\cref{sec:section43}. At \SI{400}{\milli\second}, the second line
is connected to the same busbar. The traces \texttt{Ix(Line3a:L1\_1)} and
\texttt{Ix(Line3b:L1\_1)} initially show transients of similar magnitude in both
lines.

Closer inspection shows that the high-frequency components have opposite
signs. They are capacitive currents circulating through the busbar and
the phase-to-earth capacitances of the two cables. For two identical
\SI{5}{\kilo\metre} cables, the approximate circulating-current frequency is

\begin{calculation}{Circulating-current and network-current frequencies}{calculation42}
\begin{align}
L_{\mathrm{ring}}
 &=\frac{2\cdot\SI{0.11}{\ohm\per\kilo\metre}
              \cdot\SI{2.5}{\kilo\metre}}
             {2\pi\cdot\SI{50}{\hertz}}
   \approx\SI{1.75}{\milli\henry},\\
C_{\mathrm{ring}}
 &=\frac{\SI{400}{\nano\farad}}{2}
   =\SI{200}{\nano\farad},\\
f_{\mathrm{ring}}
 &=\frac{1}{2\pi\sqrt{L_{\mathrm{ring}}C_{\mathrm{ring}}}}
   \approx\boxed{\SI{3.8}{\kilo\hertz}}.
\end{align}

The lower-frequency component drawn from the network source sees the combined
capacitance of both lines:

\begin{align}
L_{\mathrm{net}}
 &=\frac{\SI{1.613}{\ohm}
           +\SI{0.11}{\ohm\per\kilo\metre}\cdot\SI{2.5}{\kilo\metre}}
          {2\pi\cdot\SI{50}{\hertz}}
   \approx\SI{6.0}{\milli\henry},\\
C_{\mathrm{net}}&=2\cdot\SI{400}{\nano\farad}
   =\SI{800}{\nano\farad},\\
f_{\mathrm{net}}
 &=\frac{1}{2\pi\sqrt{L_{\mathrm{net}}C_{\mathrm{net}}}}
   \approx\boxed{\SI{1.03}{\kilo\hertz}}.
\end{align}
\end{calculation}

Because this network-current component has the lower frequency, its associated
inductive voltage drop and busbar-voltage transient are also smaller. The
reactive-power exchange of each cable can be viewed by holding \textsc{Alt}
and right-clicking the respective line element. LTspice then displays the
instantaneous product of element voltage and current.
The lower frequency can be observed by measuring the current at the voltage source. The higher frequency is measured by measuring the current flowing through the respective cable.

\section{Energization of a Transformer-Terminated Line}
\label{sec:section46}
NOTE: Simulate \texttt{Switching\_lines.asc} and choose after simulation \texttt{line4.plt}.
This is done by activating the waveform-viewer window and click in the PlotSettings the menu item "open plot settings file".
  
In the fourth circuit, the cable supplies a transformer initialized with
unfavourable residual flux. The transient current is therefore a combination
of capacitive cable-charging current and inductive transformer-inrush current.

In the Phase A closed at a voltage zero crossing, the capacitive contribution is
small while the inductive inrush contribution is large. In the other phases,
the capacitive component is initially dominant but decays rapidly; the
inductive component subsequently prevails. The two components decay at
different rates. Now extend the simulation to 15 s (\texttt{Switching\_lines\_15s.asc}) and run it again. During the first few seconds, the voltages on the primary side of the transformer exceed their nominal values. The Phase C voltage is particularly high. The peak value is 240kV compared to the nominal value of 90kV nominal peak voltage of line to ground voltage.
You can get this under control by operating the Petersen coil with slight undercompensation or overcompensation
(out of resonance). The residual-flux initialization and its coordination with the switching time are
described in \cref{chap:chapter3}.

\section{Extended-Line Study: Charging Current Near Rated Current}
\label{sec:section47}
\noindent NOTE: Simulate \texttt{Switching\_lines\_70km.asc} and choose after simulation for plotting \texttt{lines\_70km\_1.plt} for waveform1 or \texttt{lines\_70km\_2.plt} for waveform2

The cable length is increased until its steady-state charging current is close
to the nominal current of approximately \SI{900}{\ampere} peak, or
\SI{630}{\ampere} RMS. The simulation study uses a line length of
\SI{70}{\kilo\metre}.

For a phase-to-earth capacitance $C'$ and line-to-line RMS voltage
$U_{\mathrm{LL}}$, the approximate charging current is

\begin{equation}
I_{\mathrm c}=\omega C'l\frac{U_{\mathrm{LL}}}{\sqrt{3}}.
\label{eq:cable-charging-current}
\end{equation}

With $C'=\SI{400}{\nano\farad\per\kilo\metre}$, the ideal length required for
exactly \SI{600}{\ampere} RMS is about \SI{75}{\kilo\metre}. The
\SI{70}{\kilo\metre} simulation case is therefore a practical approximate
setting, with the exact result determined by the implemented line model.

At \SI{70}{\kilo\metre}, the current transient reaches approximately
\SI{3.5}{\kilo\ampere} in phase B and C, while its dominant oscillation frequency falls to
approximately \SI{230}{\hertz}. Despite the larger current, the voltage
transient is smaller and the oscillation decays rapidly. The greater cable
resistance provides stronger damping, and the lower oscillation frequency
reduces the network impedance.

\vspace{1cm}

\begin{calculation}{Comparison of inductive voltage drops across the network impedance}{calculation43}
For the \SI{5}{\kilo\metre} case,
\begin{align}
L_{\mathrm N}&=\frac{\SI{1.613}{\ohm}}
 {2\pi\cdot\SI{50}{\hertz}}\approx\SI{5.1}{\milli\henry},\\
X_L(\SI{1.44}{\kilo\hertz})
 &=2\pi\cdot\SI{1.44}{\kilo\hertz}\cdot L_{\mathrm N}
   \approx\SI{46.4}{\ohm},\\
\widehat U_L&\approx\SI{46.4}{\ohm}\cdot\SI{1.5}{\kilo\ampere}
  \approx\SI{69.6}{\kilo\volt}.
\end{align}

For the \SI{70}{\kilo\metre} case,
\begin{align}
X_L(\SI{230}{\hertz})
 &=2\pi\cdot\SI{230}{\hertz}\cdot L_{\mathrm N}
   \approx\SI{7.4}{\ohm},\\
\widehat U_L&\approx\SI{7.4}{\ohm}\cdot\SI{3.5}{\kilo\ampere}
  \approx\SI{25.9}{\kilo\volt}.
\end{align}
\end{calculation}

\vspace{1cm}

This explains why the \SI{70}{\kilo\metre} case produces a smaller voltage
transient even though its charging-current transient is larger. In LTspice,
the voltage drop can be measured by right-clicking the source-side terminal of the phase B impedance, choosing \textsc{Mark Reference}, and selecting the switch-side terminal of the impedance
 as the second measurement point. The observed drop is approximately
\SI{25}{\kilo\volt}, compared with approximately \SI{70}{\kilo\volt} in the
\SI{5}{\kilo\metre} case. This is because the dominant frequency is much lower and therefore the impedance is much smaller. HINT: Observe the Petersen-coil current in the \SI{70}{\kilo\metre} cable network. Its peak value is approximately \SI{90}{\ampere}.  

\newpage

\begin{figure}[H]
	\centering
	\includegraphics[width=1\textwidth] {figures/04_Switching\_lines/Switching\_lines\_schematic}
	\caption{Switching lines - schematic}
	\label{fig:figure41}
\end{figure}

\begin{figure}[H]
	\centering
	\includegraphics[width=0.9\textwidth] {figures/04_Switching\_lines/Switching\_lines\_waveforms1}
	\caption{Switching lines1 - waveform}
	\label{fig:figure42}
\end{figure}

\begin{figure}[H]
	\centering
	\includegraphics[width=0.9\textwidth] {figures/04_Switching\_lines/Switching\_lines\_waveforms2}
	\caption{Switching lines2 - waveform}
	\label{fig:figure43}
\end{figure}

\begin{figure}[H]
	\centering
	\includegraphics[width=0.9\textwidth] {figures/04_Switching\_lines/Switching\_lines\_waveforms3ab}
	\caption{Switching lines3 - waveform}
	\label{fig:figure44}
\end{figure}

\begin{figure}[H]
	\centering
	\includegraphics[width=0.9\textwidth] {figures/04_Switching\_lines/Switching\_lines\_waveforms4}
	\caption{Switching lines4 - waveform}
	\label{fig:figure45}
\end{figure}

\begin{figure}[H]
	\centering
	\includegraphics[width=0.9\textwidth] {figures/04_Switching\_lines/Switching\_lines\_15s_waveform}
	\caption{Switching lines 15s - waveform}
	\label{fig:figure48}
\end{figure}

\begin{figure}[H]
	\centering
	\includegraphics[width=1\textwidth] {figures/04_Switching\_lines/Switching\_lines\_70km_schematic}
	\caption{Switching lines 70km - schematic}
	\label{fig:figure46}
\end{figure}

\begin{figure}[H]
	\centering
	\includegraphics[width=0.9\textwidth] {figures/04_Switching\_lines/Switching\_lines\_70km_waveform1}
	\caption{Switching lines 70km - waveform1}
	\label{fig:figure47}
\end{figure}

\begin{figure}[H]
	\centering
	\includegraphics[width=0.9\textwidth] {figures/04_Switching\_lines/Switching\_lines\_70km_waveform2}
	\caption{Switching lines 70km - waveform2}
	\label{fig:figure48}
\end{figure}



\begin{chaptersummary}
Cable energization produces high-frequency oscillations determined mainly by
network inductance and cable capacitance. The switching instant controls the
initial phase currents and overvoltages, while a connected load increases
damping. Energizing a second cable creates both high-frequency circulating
currents and a lower-frequency network-current component. When a cable
supplies a transformer, capacitive charging and inductive inrush currents are
superimposed. Increasing the cable length raises the charging current but also
lowers the resonant frequency and increases resistive damping; the resulting
voltage transient can therefore become smaller even when the current transient
is larger. As can be seen in the example with the transformer, operating the Petersen coil exactly at the resonance point can cause very high voltages. For this reason, in practice, the Petersen coil is operated with either slight overcompensation or undercompensation. As can be seen in the network with the \SI{70}{\kilo\meter} cable, the residual current in very large cable networks becomes quite high. The network is therefore usually divided into smaller subnetworks.
\end{chaptersummary}
